% !TeX spellcheck = en_US
\documentclass[a4paper,12pt]{article}
\usepackage{hyperref}
\usepackage{cite}
\usepackage{graphicx}

\begin{document}

\title{Documentation pages for AWPMD package in LAMMPS}
\date{\today}
\maketitle

%===============================================================================
\section*{pair\_style wpmd/cut command}
%===============================================================================

\subsection*{Description}

This pair style represents the basic implementation of the Wave Packet Molecular Dynamics (WPMD)~\cite{Klakow:JCP94} for studying nonideal (strongly coupled) systems of charged particles such as the nonideal plasma and warm dense matter. This method is an extension of the classical molecular dynamics (MD) of electrons and ions, the ions (nuclei) are treated as classical point-like particles and the electrons are  represented as normalized Gaussian wave packets with dynamical width (size). It allows to study both equilibrium states and non-equilibrium processes beyond the Born-Oppenheimer approach due to explicit dynamics of electrons. At the moment the method is verified for hydrogen and helium plasmas although it is expected to be applicable for heavier atoms as well.

In this method, the single-electron wave function is parametrized by a set of eight scalars: the wave packet position $\mathbf{r}$ (3d vector), the wave packet width $s$ (scalar) and their conjugate momenta $\mathbf{p}$ (3d vector), $p_s$ (scalar):
\begin{equation}
\varphi(\mathbf{x}) = \left( \frac{3}{2\pi s^2} \right)^{3/4}
\exp \left\{
  - \left(\frac{3}{4s^2} - \frac{\mathrm{i}{p_s}}{2\hbar s} \right)
  (\mathbf{x}-\mathbf{r})^2 + \frac{\mathrm{i}}{\hbar}{\mathbf{p}} \cdot (\mathbf{x}-\mathbf{r})
\right\}.
\end{equation}

The electron Force Field (eFF) (see \href{pair_eff.html}{pair\_style eff/cut}) was the first pair style of such kind implemented in LAMMPS  being in fact an extension of the original WPMD algorithm where the spin-dependent Pauli potential is added (see below). The definition of the wave packet width (size) $s$ in eFF differs from this pair style by a factor of $3^{1/2}$.

Within the Hartree approximation the many-electron wave function is given as
\begin{equation}
\Psi(\{\mathbf{x}_k\}) = \prod_{k=1}^{N_\mathrm{e}}\varphi(\mathbf{x}_k),
\end{equation}
where $N_\mathrm{e}$ is the number of electrons.

The Hamiltonian has the form
\begin{equation}
H_\mathrm{wpmd}
  = \left\langle \Psi \right| \hat{H}_\mathrm{wpmd} \left| \Psi \right\rangle
  = K_\mathrm{i} + K_\mathrm{e} + K'_\mathrm{e} + U_\mathrm{ii} + U_\mathrm{ei} + U_\mathrm{ee} + U_\mathrm{ext},
\end{equation}
where $K_\mathrm{i}$ and $K_\mathrm{e} + K'_\mathrm{e}$ are the kinetic energies electrons and ions, $U_\mathrm{ii}$, $U_\mathrm{ei}$, $U_\mathrm{ee}$ are the potential energies of ion-ion, electron-ion and electron-electron interactions respectively:
\begin{eqnarray}
&& K_\mathrm{i} = \sum\limits_{k=1}^{N_\mathrm{i}} \frac{\mathbf{p_\mathrm{i}}_k^2}{2m_\mathrm{i}},
   \nonumber\\
&& K_\mathrm{e} = \sum_{k=1}^{N_\mathrm{e}} \frac{\mathbf{p}_k^2}{2m_\mathrm{e}},
   \nonumber\\
&& K'_\mathrm{e} = \sum_{k=1}^{N_\mathrm{e}} \left( \frac{9\hbar^2}{8m_\mathrm{e}s^2_k} 
+ \frac{{p_s}_k^2}{2m_\mathrm{e}} \right),
   \nonumber\\
&& U_\mathrm{ii} = \sum\limits_{k<l}^{N_\mathrm{i},\,N_\mathrm{i}}\!
   \frac{Z^2 e^2}{\left| \mathbf{R}_k - \mathbf{R}_l \right|},
   \nonumber\\
&& U_\mathrm{ei} = 
   - \sum_{k,l}^{N_\mathrm{e},\,N_\mathrm{i}}\!
   \frac{Z e^2}{|\mathbf{r}_k-\mathbf{R}_l|}\,
   \mathrm{erf} \!\Bigg(\frac{\sqrt{3}|\mathbf{r}_k - \mathbf{R}_l|}{\sqrt{2}s_k}\Bigg),
   \nonumber\\
&& U_\mathrm{ee} = \sum_{k<l}^{N_\mathrm{e},\,N_\mathrm{e}}\!
   \frac{e^2}{|\mathbf{r}_k - \mathbf{r}_l|}\,
   \mathrm{erf} \!\Bigg(\frac{\sqrt{3}|\mathbf{r}_k - \mathbf{r}_l|}{\sqrt{2}(s^2_k+s^2_l)^{1/2}}\Bigg),
\label{eq:Hwpmd}
\end{eqnarray}
$N_\mathrm{i}$ is the number of ions, $m_\mathrm{i}$ and $Ze$ are the ions mass and charge, $\mathbf{R}_k$ and ${p_\mathrm{i}}_k$ are the ions position and momentum, $m_\mathrm{e}$ and $e$ are the electron mass and charge, $U_\mathrm{ext}$ is an external potential like the wall boundary. Note that although $K'_\mathrm{e}$ is the kinetic energy, in the log and dump files it is assigned to the potential energy in order to keep the definition of kinetic energy of electron $K_\mathrm{e}$ similar to the classical MD.

The equations of motion follow from the time-dependent Schrodinger equation. In the case of Hartree approximation they correspond to the  Hamiltonian equations where the width of each electron represents an additional degree of freedom. 

The pair style has only one parameter: Rc is the cutoff radius.

The pair style is designed to be used with \href{atom_style.html}{atom\_style wavepacket} definitions, in order to handle the description of systems with interacting ions and explicit electrons.

There are a few modifications of the original WPMD algorithm to account for the antisymmetry in the many-electron wave function which is related to the exchange-correlation effects, electron degeneracy, and Pauli blocking. In all these methods the spin projections (up or down) are constantly attributed to electrons to be used in interaction model. Below is the list of such modifications included in LAMMPS:

\begin{itemize}

\item electron Force Field (eFF), see \href{pair_eff.html}{pair\_style eff/cut}. The antisymmetry is implemented via the spin-dependent semi-empirical Pauli potential. The method is almost as fast as the original WPMD.

\item Antisymmetrized Wave Packet Molecular Dynamics (AWPMD), see \href{pair_awpmd.html}{pair\_style awpmd/cut}. In AWPMD, the many-body wave function for electrons with the same spin projection is antisymmetrized using the unrestricted Hartree-Fock approximation. This pair style also supports representation of a single electron wave function by multiple Gaussian wave packets which improves the accuracy of bound state energies. Due to computations of the norm-matrix and the modified equations of motion, this method is slower than the original WPMD or eFF.

\item The joint method of the Wave Packet Molecular Dynamics and the Density Functional Theory (WPMD-DFT), see \href{pair_wpmd.html}{pair\_style wpmd/dft/cut}. In this method the exchange-correlation effects are included by and additional energy computed as a functional of the local electron density following the idea of DFT within the LSDA approximation. The electron density is obtained from the wave packet positions and widths. This method is slower than eFF but faster than AWPMD. It has a GPU-accelerated version.

\end{itemize}

The original WPMD method has a known problem of unlimited broadening of the wave packets for weakly bound electrons. Therefore the periodic boundaries are appropriate only for very high electron density. This problem can be solved either by manual limiting of the wave packet width (eFF) or by using a 3-dimensional confining potential (wall potential) for the whole system which naturally constrains both the wave packet positions and widths (see \href{fix_wall_wpmd.html}{fix wall/wpmd}).

The time-dependent dynamics of wave packets is implemented using \href{fix_nve_wpmd.html}{fix nve/wpmd} for the energy conservative system and  \href{fix_nh_wpmd.html}{fix nvt/wpmd, npt/wpmd, nph/wpmd} for the canonical, isothermal-isobaric, and isenthalpic ensembles. The Monte-Carlo sampling is given by \href{fix_mc_wpmd.html}{fix mc/wpmd}. It can be used for all mentioned WPMD modifications.


%-------------------------------------
\bibliographystyle{ieeetr}
\bibliography{morozov,nplasma,computing}

\end{document}
